\section{Covering bounds for scalar separators}
\label{sec:scalar-covering}

This appendix bounds the number of separator cells in the exact-bag
split relaxation of \cref{def:consistency-split}.  It does not count local
leaves or replace the bag relaxations in a decomposition certificate.
We use the value functions $U_t,V_t,L_t$, the separator margin
$w_t=U_t-L_t$, and the graded split
$\psi_t=U_t-\theta_t w_t$ from
\cref{eq:consistency-band,thm:consistency-graded}.

\subsection{Concentration and affine approximation}

A finite-valued function $U$ on an interval is $M$-semiconcave if
$U(s)-Ms^2/2$ is concave.  A function $L$ is $M$-semiconvex if
$L(s)+Ms^2/2$ is convex.  For a nondegenerate closed interval
$J=[a,b]$ with finite one-sided endpoint derivatives, put
\begin{align}
 \Delta_U(J)&=U'(a+)-U'(b-)+M(b-a),\\
 \Delta_L(J)&=L'(b-)-L'(a+)+M(b-a).
 \label{eq:scalar-covering-slope-mass}
\end{align}
These quantities are nonnegative and increase when $J$ is enlarged.
They are the slope drop of $U(s)-Ms^2/2$ and the slope rise of
$L(s)+Ms^2/2$, respectively.

\begin{lemma}[Scalar concentration]
\label{lem:scalar-covering-concentration}
Let $M\geq0$, let $U$ be $M$-semiconcave and $L$ be
$M$-semiconvex on an interval $I$, and suppose $L\leq U$.
Write $w=U-L$.  If $h>0$ and $[c-h,c+h]\subset I$, then
\begin{equation}
 \Delta_U([c-h/2,c+h/2])+
 \Delta_L([c-h/2,c+h/2])
 \leq \frac{4w(c)}h+4Mh.
 \label{eq:scalar-covering-concentration}
\end{equation}
\end{lemma}

\begin{proof}
Set
\[
 U_c(s)=U(s)-\frac M2(s-c)^2,
 \qquad L_v(s)=L(s)+\frac M2(s-c)^2,
 \qquad \varphi=L_v-U_c.
\]
The function $\varphi$ is convex, satisfies
$\varphi(s)\leq M(s-c)^2$, and has $\varphi(c)=-w(c)$.
For $a=c-h/2$ and $b=c+h/2$, the left side of
\eqref{eq:scalar-covering-concentration} is
$\varphi'(b-)-\varphi'(a+)$.  Convexity gives
\[
 \varphi'(b-)\leq\frac{2}{h}\bigl(\varphi(c+h)-\varphi(b)\bigr),
 \qquad
 \varphi'(a+)\geq\frac{2}{h}\bigl(\varphi(a)-\varphi(c-h)\bigr).
\]
Also $\varphi(a)+\varphi(b)\geq2\varphi(c)$.  Therefore
\begin{align*}
 \varphi'(b-)-\varphi'(a+)
 &\leq\frac2h\bigl(\varphi(c+h)+\varphi(c-h)-2\varphi(c)\bigr)\\
 &\leq\frac2h\bigl(2Mh^2+2w(c)\bigr).
\end{align*}
The one-sided derivatives used here are finite because $a$ and $b$ lie
in the interior of $[c-h,c+h]$.
\end{proof}

\begin{lemma}[Chord oscillation]
\label{lem:scalar-covering-chord}
Let $D=[c-r,c+r]$ with $r>0$.  Suppose $U$ is
$M$-semiconcave, $L$ is $M$-semiconvex, and their one-sided endpoint
derivatives on $D$ are finite.  For every $\theta\in[0,1]$, the function
$\psi=(1-\theta)U+\theta L$ has an affine approximation $l$ satisfying
\begin{equation}
 \osc_D(\psi-l)
 \leq\frac r2\bigl(\Delta_U(D)+\Delta_L(D)\bigr)
       +\frac M2r^2.
 \label{eq:scalar-covering-oscillation}
\end{equation}
\end{lemma}

\begin{proof}
First let $f$ be concave on $[a,b]$, with finite endpoint slopes
$p=f'(a+)$ and $q=f'(b-)$, and let $l_f$ be its chord with slope $\mu$.
Then $q\leq\mu\leq p$, and the endpoint supporting lines give
\[
 0\leq f(x)-l_f(x)
 \leq\min\{(p-\mu)(x-a),(\mu-q)(b-x)\}.
\]
If $p=q$, then $f=l_f$.  Otherwise the maximum of the right side is
\[
 (b-a)\frac{(p-\mu)(\mu-q)}{p-q}
 \leq\frac{b-a}{4}(p-q).
\]
Applying the same argument to $-f$ proves the corresponding bound for
the distance of a convex function below its chord.

Now $U_c=U-M(s-c)^2/2$ is concave and
$L_v=L+M(s-c)^2/2$ is convex.  Let $l_U,l_L$ be their chords on $D$,
and take
\[
 l=(1-\theta)l_U+\theta l_L
      +(1-2\theta)\frac M2r^2.
\]
The residuals $U_c-l_U$ and $L_v-l_L$ have oscillations at most
$r\Delta_U(D)/2$ and $r\Delta_L(D)/2$, respectively.
The remaining quadratic residual is
$(1-2\theta)M((s-c)^2-r^2)/2$, whose oscillation is at most $Mr^2/2$.
Oscillation is subadditive, so
\[
 \osc_D(\psi-l)
 \leq\frac r2\bigl((1-\theta)\Delta_U(D)
                      +\theta\Delta_L(D)\bigr)+\frac M2r^2,
\]
which implies \eqref{eq:scalar-covering-oscillation}.
\end{proof}

\subsection{The scalar covering theorem}

Assume there are $N$ bags and $n_T=N-1\geq1$ edges.  Identify each edge
with its nonroot bag $t$.  Suppose every separator is one coordinate
with nondegenerate compact range $I_t$ of length $\ell_t>0$.
The bag functions are continuous on their compact boxes.  For each
$t\ne r$, choose strictly positive constants $M_t,G_t$ such that $U_t$
is $M_t$-semiconcave, $L_t$ is $M_t$-semiconvex, and both are
$G_t$-Lipschitz on $I_t$.  A zero regularity constant may always be
replaced by a positive upper bound.

\begin{proposition}[Regularity from smooth bag data]
\label{prop:scalar-covering-smooth}
If each bag objective $a_u$ is $C^1$ on a neighborhood of its compact
product box and has a Lipschitz gradient there, then the preceding
semiconcavity, semiconvexity, and Lipschitz assumptions hold with finite
positive constants.
\end{proposition}

\begin{proof}
Fix a scalar separator $S_t$ and write its coordinate as $s$.
The objective on either side of the edge is a sum of bag objectives.
For every fixed choice of the other coordinates, this sum has a
derivative in $s$ whose Lipschitz constant is bounded uniformly by the
sum of the bag gradient Lipschitz constants on that side.
Its derivative is also uniformly bounded, because the bag gradients
are continuous on compact boxes.
Thus each side objective is uniformly semiconcave and uniformly
Lipschitz in $s$, with finite constants.

The remaining coordinates range over a product box independent of
$s$.  If $H(s,y)$ is one of the side objectives and $M$ is its uniform
semiconcavity constant, then
\[
 \inf_y H(s,y)-\frac M2s^2
 =\inf_y\left(H(s,y)-\frac M2s^2\right)
\]
is concave: a pointwise infimum of concave functions is concave.
The same infimum is Lipschitz with the uniform Lipschitz constant,
since $H(s,y)\leq H(s',y)+G\abs{s-s'}$ for every $y$,
and one can take infima and then interchange $s,s'$.
Applying these observations to the child side gives semiconcavity and
Lipschitz continuity of $U_t$, and applying them to the other side
gives the same properties for $V_t$.
Consequently $L_t=f^*-V_t$ is semiconvex and Lipschitz.
Taking positive upper bounds for the two sides supplies $M_t,G_t$.
\end{proof}

Fix $\eps>0$ and $\tau=1/(2n_T)$.  For a partition $P_t$ into intervals,
let $\operatorname{PA}(P_t)$ be the class of independent affine pieces
on its cells, with a fixed convention assigning every shared endpoint
to one cell.  Write
$\Phi(P)=\prod_{t\ne r}\operatorname{PA}(P_t)$.
All cell suprema below are taken on the closed cell, so they also bound
this endpoint convention.  Define the closed-cell bracket
\begin{equation}
 g_{t,D}=\inf_{l\text{ affine}}
 \left\{\sup_D(l-\psi_t-\tau w_t)
             +\sup_D(\psi_t-\tau w_t-l)\right\}.
 \label{eq:scalar-covering-bracket}
\end{equation}
The cell sandwich \eqref{eq:consistency-cell-sandwich} gives
$\operatorname{gap}(\Phi(P))\leq
\sum_{t\ne r}\max_{D\in P_t}g_{t,D}$.

Starting from $I_t$, bisect a cell $D$ of radius $r_D$ while
$B_t(D)>\eps/n_T$, where
\begin{equation}
 B_t(D)=
 \begin{cases}
  (8n_T+\tfrac12)M_t r_D^2-\dfrac{\min_D w_t}{2n_T},
     &\dist(D,\partial I_t)\geq4n_T r_D,\\[2mm]
  2G_t r_D+\tfrac52M_t r_D^2-\dfrac{\min_D w_t}{n_T},
     &\dist(D,\partial I_t)<4n_T r_D.
 \end{cases}
 \label{eq:scalar-covering-rule}
\end{equation}
On each final cell use a minimizer of
\eqref{eq:scalar-covering-bracket}, with its intercept shifted so that
the two suprema are equal.

For $A\subset\R$ and $d>0$, let $N_\infty(A,d)$ be the least number of
closed intervals of length $d$ covering $A$, with
$N_\infty(\varnothing,d)=0$.  Put
\begin{align}
 E(\eta)&=\{x\in X:F(x)-f^*\leq\eta\},\\
 \mathcal N_t&=\sup_{\eta\geq0}
 N_\infty\!\left(\pi_{S_t}E(\eta),
                   2\sqrt{\frac{\eps+\eta}{M_t}}\right),\\
 J_t&=\max\left\{0,
 \ceil{\log_2\left(\ell_t\sqrt{
             \frac{(16n_T^2+n_T)M_t}{8\eps}}\right)}\right\},\\
 \varrho_t&=\min\left\{\frac{\eps}{4n_T G_t},
                         \sqrt{\frac{\eps}{5n_T M_t}}\right\},\\
 J'_t&=\max\left\{0,\ceil{\log_2
                      \left(\frac{\ell_t}{2\varrho_t}\right)}\right\}.
 \label{eq:scalar-covering-numbers}
\end{align}
The covering number $\mathcal N_t$ is finite, since all covering
intervals have length at least $2\sqrt{\eps/M_t}$; it is at least one
because $E(0)$ is nonempty.

\begin{theorem}[Covering bound for scalar separators]
\label{thm:scalar-covering}
Under the preceding assumptions, the rule
\eqref{eq:scalar-covering-rule} terminates and produces an exact-bag
split with gap at most $\eps$.  Its separator partitions satisfy
\begin{equation}
 |P_t|\leq1+(4n_T+2)J_t\mathcal N_t+4n_T J'_t
 \qquad(t\ne r).
 \label{eq:scalar-covering-count}
\end{equation}
In particular, $\operatorname{gap}(\Phi(P))\leq\eps$.
\end{theorem}

\begin{proof}
We first prove $g_{t,D}\leq B_t(D)$ for every dyadic cell.
For a cell in the first case of \eqref{eq:scalar-covering-rule}, choose
$c_m\in D$ minimizing $w_t$ and put $h=4n_T r_D$.
Then $[c_m-h,c_m+h]\subset I_t$ and
$D\subset[c_m-h/2,c_m+h/2]$.  Monotonicity of the slope masses and
\cref{lem:scalar-covering-concentration,lem:scalar-covering-chord}
give an affine $l$ with
\[
 \osc_D(\psi_t-l)
 \leq\frac{\min_D w_t}{2n_T}
       +(8n_T+\tfrac12)M_t r_D^2.
\]
For every affine $l$, the bracket in
\eqref{eq:scalar-covering-bracket} is at most
$\osc_D(\psi_t-l)-2\tau\min_D w_t$.  This proves the required bound
for the first case.

For any cell, Lipschitz continuity bounds every one-sided derivative of
$U_t$ and $L_t$ by $G_t$ in absolute value.  Consequently
\[
 \Delta_{U_t}(D)+\Delta_{L_t}(D)\leq4G_t+4M_t r_D.
\]
\Cref{lem:scalar-covering-chord} now bounds the oscillation by
$2G_t r_D+\tfrac52M_t r_D^2$, proving the second case.
Since $w_t\geq0$, both upper bounds tend uniformly to zero as the
cell radius tends to zero.  Thus the rule terminates after finitely
many levels, and every final cell has $g_{t,D}\leq\eps/n_T$.

The affine minimizers used by the rule exist.  Indeed, after writing
$l(s)=as+b$, the intercept cancels in the bracket.  The resulting
function of $a$ is continuous.  Evaluating its two suprema at opposite
endpoints of $D$ bounds it below by $2r_D\abs{a}$ minus a fixed
constant, so it is coercive and attains its minimum.  Shifting $b$
balances the two suprema.  The proof of the cell sandwich then bounds
the gap of the resulting split by
$\sum_{t\ne r}\max_{D\in P_t}g_{t,D}\leq\eps$.

It remains to count the split cells.  Fix $t$ and abbreviate
$M=M_t$, $G=G_t$, and $\ell=\ell_t$.
At level $j\geq0$, all dyadic cells have radius
$r_j=\ell\,2^{-j-1}$.
A split cell in the first case must satisfy
\[
 \min_D w_t<\eta_j:=(16n_T^2+n_T)M r_j^2-2\eps.
\]
In particular $\eta_j>0$.  Such a cell meets
$\{s:w_t(s)\leq\eta_j\}=\pi_{S_t}E(\eta_j)$; the equality follows
from the attained separator-margin minimum on a compact box.
There are at most $J_t$ levels with $\eta_j>0$, since this condition is
equivalent to
\[
 2^j<\ell\sqrt{\frac{(16n_T^2+n_T)M}{8\eps}}.
\]
At each such level, cover $\pi_{S_t}E(\eta_j)$ by at most
$\mathcal N_t$ intervals of length
$d_j=2\sqrt{(\eps+\eta_j)/M}$.  Their length relative to a cell is
bounded by
\[
 \frac{d_j}{2r_j}
 =\sqrt{16n_T^2+n_T-\frac{\eps}{Mr_j^2}}
 <\sqrt{16n_T^2+n_T}<4n_T+1.
\]
A closed interval of length $d_j$ meets at most
$\floor{d_j/(2r_j)}+2\leq4n_T+2$ closed dyadic cells.
There are therefore at most $(4n_T+2)J_t\mathcal N_t$ split cells in
the first case.

At any level, fewer than $2n_T$ cell lengths separate a cell in the
second case from one of the two endpoints of $I_t$.
There are at most $2n_T$ such cells at each end, and hence at most $4n_T$
in total.  A cell in this case can split only if
\[
 2G r_j+\tfrac52 M r_j^2>\eps/n_T.
\]
If $r_j\leq\varrho_t$, each term on the left is at most
$\eps/(2n_T)$, so the cell cannot split.
The condition $r_j>\varrho_t$ holds at at most $J'_t$ levels.
Thus the second case contributes at most $4n_T J'_t$ split cells.
Each bisection increases the final cell count by one, which proves
\eqref{eq:scalar-covering-count}.
\end{proof}

\begin{corollary}[Refinement by the cell bracket]
\label{cor:scalar-covering-bracket}
Under the assumptions of \cref{thm:scalar-covering}, start from $I_t$
and bisect a dyadic cell $D$ while $g_{t,D}>\eps/n_T$.
Use balanced bracket minimizers on the final cells.
This rule terminates, gives an exact-bag split with gap at most $\eps$,
and produces a coarsening of the partition in
\cref{thm:scalar-covering}.  Its cell count therefore satisfies
\eqref{eq:scalar-covering-count}.
\end{corollary}

\begin{proof}
The proof of \cref{thm:scalar-covering} establishes
$g_{t,D}\leq B_t(D)$ for every dyadic cell.  Hence any cell split by the
bracket rule is also split by \eqref{eq:scalar-covering-rule}.
Induction on the level shows that every ancestor needed to reach such
a cell is split by both rules.  The bracket rule therefore uses a
subset of the finitely many bisections in
\cref{thm:scalar-covering}; its final partition is a coarsening.
Its final brackets are at most $\eps/n_T$, so the cell sandwich gives
the asserted gap bound.
\end{proof}
